\section{敏感性与固定方案压力边界}\label{sec:sensitivity}
四问得到的是官方确定性参数下的结果。为判断其可解释范围，须区分两种问题：改变参数后允许重新组批，最优目标怎样变化；保持现有计划不动，还能承受多少额外压力。前者反映重新决策后的适应性，后者只反映现有方案的余量，二者不能互相替代。本文只使用已计算的确定性场景，不构造未经观测的概率分布。

\subsection{储备变化与离散组批跳变}
储备提高直接收紧式\eqref{eq:feasible}，而不改变额定载荷。对 10\%--40\% 储备取 \ReserveGridCount{} 个等距点，重算服务区—机型安全载荷，共 \CapacityCaseCount{} 个能力场景；另在表\ref{tab:sensitivity}所列的 \BatchScenarioCount{} 个储备场景重新运行 Q1 精确组批。前者显示连续能力收缩，后者提供基准交叉核验。另枚举全部质量与体积可行子集，取各子集在官方机型中的最大可行储备，再用“最大化最小储备”的集合划分DP求每个精确架次数的储备上限；按所有临界点划分参数区间，并在各区间内部用原Q1求解器独立核对最少架次。因此图中的阶梯不是稀疏场景插值。
\paperfig{fig11_sensitivity}{提高储备使连续运载能力收缩，重新组批的离散架次数随之改变。上图细线为全部服务区，粗线追踪基准最受限服务区，不可达用轴下叉号单独标识；下图阶梯来自完整子集临界分析，空心点为原五场景独立重算，浅色斜纹区表示不可行；灰虚线分别为20\%基准和固定原方案储备上限。}{sensitivity}
\papertable{sensitivity}{储备增大对能力与组批的影响不同：表中每个场景均重新精确求解}{sensitivity}
基准及两个较低储备场景的组批目标保持一致，而更高储备下需要增加架次。原因不是算法突然失效，而是货箱不可连续切小：原先一个批次的固定质量越过能力边界时，必须换机型、拆成新的货箱子集或增加架次。因此，连续的物理参数变化会产生离散的组合决策跳变。

基准固定组批最低 SOC 为 \QOneReserveLimit\%，故不改变货箱、路线和机型时，最大共同储备就是该值。单凭这一固定方案值不能判断所有同架次数方案。本次额外的完整子集临界分析独立得到：最少架次首次在约 \ReserveTripBreak\% 后增加，该临界值在本算例恰与固定基准阈值重合；储备超过约 \ReserveReachabilityLimit\% 后，至少一个服务区不存在可行货箱划分。阶梯在临界点本身仍取较低架次数，图中端点和数值容差沿用统一能源可行判据。

\subsection{固定运输计划的能耗容忍度}
设所有运输架次的计算能耗统一乘以 $\lambda$，同时冻结路径、载荷与时间，则维持返航储备的最大系数为
\begin{equation}
\lambda_{\max}=\min_{p\in\mathcal P}
\frac{(1-\rho_{g_p})E_{g_p}}{E_p}.\label{eq:stress}
\end{equation}
每个架次分别给出一个上界，取最小值即可，无需新增随机仿真。Q2 与 Q3 当前运输计划的能耗上浮容忍度分别约为 \QTwoEnergyTolerance\% 与 \QThreeEnergyTolerance\%；Q3 最低 SOC 为 \MinSoc\%，仅略高于题设储备。因而“全部架次满足能源约束”只能证明基准可行，不能推断能耗估计有较大误差时仍可靠。

式\eqref{eq:stress}仅检查运输返航电量。它不代表整个联合方案可无条件承受同样增幅：能耗增加还可能延长电池充电，压缩后续任务可用时间；中继能源也需另行检查。本文没有把能耗放大后重新求解全部排程，因此不报告这种更强的恢复能力。

\subsection{通信余量与可解释范围}
保持链路几何、窗口和保障关系不变，若所有链路损耗统一增加同一常数，则在既有采样集合中可承受的增量上界等于最小可用余量。本次修复后的0.5 s回放最小值为 \MarginHalf{} dB，因而仅在固定采样集合和固定几何下允许同量级的共同额外损耗。正余量很小，不能据此推断实际衰落分布或可靠率。

新的地形遮挡可能改变链路可用模式，设备误差、天气和轨迹误差也不必等价于统一常数损耗。因此，此分析不能作为连续时间或新场景保证。若继续改进，应首先关注最紧通信区间和高能耗架次，再评估调整中继、分组和排程后的资源代价；本轮已修复已发现的空间链路反例，尚未进行上述误差情景下的全面重调度。
\FloatBarrier
